\chapter{Complexity — Reductions}
%%% generated by the blueprint pipeline from TCSlib.Complexity.ClassNP.Reductions
%%%%%%%%%%%%%%%%%%%%%%%%%%%%%%%%%%%%%%%%%%%%%%%%%%%%%%%%%%%%%%%%%%%%%%%%%%%%%%%
\section{Overview}

This module introduces polynomial-time Karp reducibility $L \le_p L'$,
$\mathsf{NP}$-hardness and $\mathsf{NP}$-completeness following Arora--Barak
Definition~2.7, and proves the basic laws of Theorem~2.8: reflexivity and
transitivity of $\le_p$, downward closure of $\mathsf{P}$ under $\le_p$, and the
collapse $\mathsf{P} = \mathsf{NP}$ whenever an $\mathsf{NP}$-hard language lies
in $\mathsf{P}$. It closes with Arora--Barak Exercise~2.8: the halting language
is $\mathsf{NP}$-hard for every representation scheme of machines, yet (for
effective schemes) not in $\mathsf{NP}$, hence not $\mathsf{NP}$-complete; the
supporting material builds a halting recognizer from a Boolean decider and a
machine that prepends a fixed word to its input.

%%%%%%%%%%%%%%%%%%%%%%%%%%%%%%%%%%%%%%%%%%%%%%%%%%%%%%%%%%%%%%%%%%%%%%%%%%%%%%%
\section{Declarations}

\begin{definition}[Polynomial-time Karp reducibility]
\lean{Complexity.PolyTimeReducible}
\label{Complexity.PolyTimeReducible}
\leanok
\uses{Complexity.PolyTimeComputable}
For languages $L, L' \subseteq \{0,1\}^{*}$, we say $L$ is polynomial-time
(Karp) reducible to $L'$, written $L \le_p L'$, if there is a function
$f : \{0,1\}^{*} \to \{0,1\}^{*}$ that is polynomial-time computable (some
finite-state binary-alphabet Turing machine computes it within $C \cdot (n+1)^c$
steps on inputs of length $n$, for constants $C, c \in \mathbb{N}$) such that for
every binary string $x$, $x \in L$ if and only if $f(x) \in L'$
[AB09, Definition~2.7].
\end{definition}

\begin{theorem}[Reflexivity of Karp reducibility]
\lean{Complexity.PolyTimeReducible.refl}
\label{Complexity.PolyTimeReducible.refl}
\leanok
\uses{Complexity.polyTimeComputable_id}
\difficulty{1}
Every language $L \subseteq \{0,1\}^{*}$ is polynomial-time Karp reducible to
itself: $L \le_p L$ [AB09, Exercise~2.9].
\end{theorem}

\begin{theorem}[Transitivity of Karp reducibility]
\lean{Complexity.PolyTimeReducible.trans}
\label{Complexity.PolyTimeReducible.trans}
\leanok
\uses{Complexity.PolyTimeComputable.comp}
\difficulty{2}
Let $L, L', L'' \subseteq \{0,1\}^{*}$ be languages. If $L \le_p L'$ and
$L' \le_p L''$, where $\le_p$ denotes polynomial-time Karp reducibility, then
$L \le_p L''$ [AB09, Theorem~2.8.1].
\end{theorem}

\begin{theorem}[$\mathsf{P}$ is closed under Karp reductions]
\lean{Complexity.mem_P_of_polyTimeReducible}
\label{Complexity.mem_P_of_polyTimeReducible}
\leanok
\uses{Complexity.P, Complexity.PolyTimeComputable, Complexity.PolyTimeComputable.comp, Complexity.mem_P_iff, Complexity.mem_P_of_dtime_le, Complexity.succ_pow_le, Turing.FinTM.DecidesInTime, Turing.MultiTapeTM.indicator}
\difficulty{4}
Let $L, L' \subseteq \{0,1\}^{*}$ be languages such that $L$ is polynomial-time
Karp reducible to $L'$. If $L' \in \mathsf{P}$, then $L \in \mathsf{P}$
[AB09, remark after Definition~2.7].
\end{theorem}

\begin{definition}[$\mathsf{NP}$-hardness]
\lean{Complexity.NPHard}
\label{Complexity.NPHard}
\leanok
\uses{Complexity.NP}
A language $L \subseteq \{0,1\}^{*}$ is $\mathsf{NP}$-hard if every language
$L' \in \mathsf{NP}$ is polynomial-time Karp reducible to it, i.e.\
$L' \le_p L$ [AB09, Definition~2.7].
\end{definition}

\begin{definition}[$\mathsf{NP}$-completeness]
\lean{Complexity.NPComplete}
\label{Complexity.NPComplete}
\leanok
\uses{Complexity.NP, Complexity.NPHard}
A language $L \subseteq \{0,1\}^{*}$ is $\mathsf{NP}$-complete if
$L \in \mathsf{NP}$ and $L$ is $\mathsf{NP}$-hard, i.e.\ every language in
$\mathsf{NP}$ is polynomial-time Karp reducible to $L$ [AB09, Definition~2.7].
\end{definition}

\begin{theorem}[An $\mathsf{NP}$-hard language in $\mathsf{P}$ collapses $\mathsf{NP}$]
\lean{Complexity.P_eq_NP_of_NPHard_mem_P}
\label{Complexity.P_eq_NP_of_NPHard_mem_P}
\leanok
\uses{Complexity.NP, Complexity.NPHard, Complexity.P, Complexity.P_subset_NP, Complexity.mem_P_of_polyTimeReducible}
\difficulty{2}
Let $L \subseteq \{0,1\}^{*}$ be an $\mathsf{NP}$-hard language, i.e.\ every
language in $\mathsf{NP}$ is polynomial-time Karp reducible to $L$. If
$L \in \mathsf{P}$, then $\mathsf{P} = \mathsf{NP}$ [AB09, Theorem~2.8.2].
\end{theorem}

\begin{theorem}[An $\mathsf{NP}$-complete language is in $\mathsf{P}$ iff $\mathsf{P}=\mathsf{NP}$]
\lean{Complexity.NPComplete.mem_P_iff}
\label{Complexity.NPComplete.mem_P_iff}
\leanok
\uses{Complexity.NP, Complexity.NPComplete, Complexity.P, Complexity.P_eq_NP_of_NPHard_mem_P}
\difficulty{2}
Let $L \subseteq \{0,1\}^{*}$ be an $\mathsf{NP}$-complete language, i.e.\
$L \in \mathsf{NP}$ and every language in $\mathsf{NP}$ is polynomial-time Karp
reducible to $L$. Then $L \in \mathsf{P}$ if and only if
$\mathsf{P} = \mathsf{NP}$ [AB09, Theorem~2.8.3].
\end{theorem}

\begin{definition}[State encoding of the halting recognizer]
\lean{Complexity.acceptState}
\label{Complexity.acceptState}
\leanok
Let $Q$ be a set of control states. The recognizer's control states are the
pairs $(q, b) \in Q \times \{0,1\}$ together with a distinguished loop state
$\ell$, and a further symbol $\bot$ marks halting. Given $s \in Q \cup \{\bot\}$
(a state, or $\bot$ for halting) and a bit $b$, the encoded state
$\mathrm{enc}(s, b)$ is $(s, b)$ when $s \in Q$, is $\bot$ when $s = \bot$ and
$b = 1$, and is the loop state $\ell$ when $s = \bot$ and $b = 0$.
\end{definition}

\begin{definition}[Recorded-bit transform of a machine action]
\lean{Complexity.acceptAction}
\label{Complexity.acceptAction}
\leanok
\uses{Complexity.acceptState, Turing.Action}
Let $a$ be a one-step action of a $k$-work-tape binary-alphabet Turing machine
with state set $Q$: an input-head move, a write-and-move instruction for each
work tape, an optional output bit, and a successor state $s \in Q \cup \{\bot\}$
($\bot$ meaning halt). For a bit $b$, the transformed action has the same
input-head move and work-tape instructions, emits no output, and has successor
state $\mathrm{enc}(s, b')$, where $b'$ is the bit emitted by $a$ (or $b$ if $a$
emits nothing) and $\mathrm{enc}(s, b')$ is $(s, b')$ for $s \in Q$, halt for
$s = \bot$ and $b' = 1$, and a dedicated loop state for $s = \bot$ and $b' = 0$.
\end{definition}

\begin{definition}[Halting recognizer of a Boolean decider]
\lean{Complexity.acceptTM}
\label{Complexity.acceptTM}
\leanok
\uses{Complexity.acceptAction, Turing.FinTM}
Let $M$ be a finite-state binary-alphabet Turing machine with $k$ work tapes,
state set $Q$ and start state $q_0$. Its halting recognizer $M^{\mathrm{acc}}$
has the same $k$ work tapes, states $(Q \times \{0,1\}) \cup \{\ell\}$ and start
state $(q_0, 0)$. In the loop state $\ell$ it moves no head, writes and emits
nothing, and stays in $\ell$; in state $(q, b)$ it performs $M$'s transition from
$q$ on the scanned symbols, except that it emits nothing, records in its bit the
symbol $M$ would emit (keeping $b$ if none), and where $M$ would halt it halts if
the recorded bit is $1$ and enters $\ell$ if it is $0$.
\end{definition}

\begin{definition}[Configuration correspondence for the halting recognizer]
\lean{Complexity.acceptCfg}
\label{Complexity.acceptCfg}
\leanok
\uses{Complexity.acceptState, Complexity.acceptTM, Turing.Cfg, Turing.FinTM}
Let $M$ be a finite-state binary-alphabet Turing machine with state set $Q$, and
$M^{\mathrm{acc}}$ its halting recognizer (states $(q,b) \in Q \times \{0,1\}$
plus a loop state $\ell$, the bit recording $M$'s last emitted symbol). For an
input $x$ and a configuration $C$ of $M$ on $x$, the configuration $\Phi(C)$ of
$M^{\mathrm{acc}}$ on $x$ has the same input-head position, work-tape contents
and work-head positions as $C$, an empty output tape, and state
$\mathrm{enc}(s, b)$, where $s$ is the state of $C$ (or $\bot$ if halted), $b$ is
the last bit on $C$'s output tape ($0$ if it is empty), and $\mathrm{enc}(s,b)$
is $(s,b)$ for $s \in Q$, halted for $s=\bot, b=1$, and $\ell$ for
$s=\bot, b=0$.
\end{definition}

\begin{lemma}[The loop state of the halting recognizer is stationary]
\lean{Complexity.acceptTM_loop}
\label{Complexity.acceptTM_loop}
\leanok
\uses{Complexity.acceptTM, Turing.Action.apply, Turing.Cfg, Turing.FinTM, Turing.MultiTapeTM.runFrom, Turing.MultiTapeTM.runFrom_succ_eq_step', Turing.MultiTapeTM.step}
\difficulty{3}
Let $M$ be a finite-state binary-alphabet Turing machine and $M^{\mathrm{acc}}$
its halting recognizer, a machine with a distinguished loop state $\ell$ in
which it moves no head, writes and emits nothing, and remains in $\ell$. For
every input $x$, every configuration $C$ of $M^{\mathrm{acc}}$ on $x$ whose state
is $\ell$, and every $t \in \mathbb{N}$, running $M^{\mathrm{acc}}$ for $t$ steps
from $C$ yields $C$ itself; in particular it never halts.
\end{lemma}

\begin{lemma}[Recorded-bit actions commute with the configuration correspondence]
\lean{Complexity.acceptCfg_apply}
\label{Complexity.acceptCfg_apply}
\leanok
\uses{Complexity.acceptAction, Complexity.acceptCfg, Turing.Action, Turing.Action.apply, Turing.Cfg, Turing.FinTM}
\difficulty{3}
Let $M$ be a finite-state binary-alphabet Turing machine with state set $Q$, $x$
an input, $C$ a configuration of $M$ on $x$, and $a$ an action of $M$. Let $b$
be the last bit on $C$'s output tape ($0$ if empty), and let $\Phi(C)$ be the
recognizer configuration with $C$'s head positions and tape contents, empty
output, and state $\mathrm{enc}(s,b)$, where $s$ is $C$'s state and
$\mathrm{enc}(s,b) = (s,b)$ for $s \in Q$, halt for a halted $s$ with $b=1$, and
a loop state for a halted $s$ with $b=0$. Then applying to $\Phi(C)$ the action
with $a$'s head moves and writes, no output, and successor state
$\mathrm{enc}(s', b')$ ($s'$ the successor of $a$, $b'$ the bit $a$ emits, or $b$
if none) gives exactly $\Phi(a(C))$, where $a(C)$ is the result of applying $a$
to $C$.
\end{lemma}

\begin{lemma}[One step of the halting recognizer simulates one step]
\lean{Complexity.acceptCfg_step}
\label{Complexity.acceptCfg_step}
\leanok
\uses{Complexity.acceptAction, Complexity.acceptCfg, Complexity.acceptCfg_apply, Complexity.acceptState, Complexity.acceptTM, Complexity.acceptTM_loop, Turing.Action.apply, Turing.Cfg, Turing.Cfg.inputSymbol, Turing.Cfg.workTapeSymbols, Turing.FinTM, Turing.MultiTapeTM.step, Turing.MultiTapeTM.step_of_halt}
\difficulty{4}
Let $M$ be a finite-state binary-alphabet Turing machine with state set $Q$, and
let $M^{\mathrm{acc}}$ be its halting recognizer: same work tapes, states
$(q,b) \in Q \times \{0,1\}$ plus a stationary loop state $\ell$, simulating $M$
with output suppressed and the last emitted bit recorded, halting where $M$
halts with recorded bit $1$ and entering $\ell$ where it halts with bit $0$. For
a configuration $C$ of $M$ on an input $x$, let $\Phi(C)$ be the configuration of
$M^{\mathrm{acc}}$ with $C$'s head positions and work-tape contents, empty output,
and state $(q,b)$, halted, or $\ell$ according as $C$ is in state $q$, halted
with last output bit $1$, or halted with last output bit $0$ (the bit $b$ being
the last output bit, $0$ if none). Then for every such $C$, one step of
$M^{\mathrm{acc}}$ from $\Phi(C)$ yields $\Phi(C')$, where $C'$ is the result of
one step of $M$ from $C$ (a halted configuration stepping to itself).
\end{lemma}

\begin{lemma}[Runs of the halting recognizer simulate runs]
\lean{Complexity.acceptTM_run}
\label{Complexity.acceptTM_run}
\leanok
\uses{Complexity.acceptCfg, Complexity.acceptCfg_step, Complexity.acceptTM, Turing.FinTM, Turing.MultiTapeTM.initCfg, Turing.MultiTapeTM.runFrom, Turing.MultiTapeTM.runFrom_comm_of_step}
\difficulty{2}
Let $M$ be a finite-state binary-alphabet Turing machine with state set $Q$, and
let $M^{\mathrm{acc}}$ be its halting recognizer: same work tapes, states
$(q,b) \in Q \times \{0,1\}$ plus a stationary loop state $\ell$, start state
$(q_0, 0)$, simulating $M$ with output suppressed and the last emitted bit
recorded, halting where $M$ halts with recorded bit $1$ and entering $\ell$ where
it halts with bit $0$. For every input $x$ and every $t \in \mathbb{N}$, the
configuration of $M^{\mathrm{acc}}$ after $t$ steps on $x$ has the same head
positions and work-tape contents as the configuration $C_t$ of $M$ after $t$
steps on $x$, an empty output tape, and state $(q,b)$, halted, or $\ell$
according as $C_t$ is in state $q$, halted with last output bit $1$, or halted
with last output bit $0$, where $b$ is the last output bit of $C_t$ ($0$ if
none).
\end{lemma}

\begin{lemma}[The halting recognizer halts exactly on accepted inputs]
\lean{Complexity.acceptTM_halts_iff}
\label{Complexity.acceptTM_halts_iff}
\leanok
\uses{Complexity.acceptState, Complexity.acceptTM, Complexity.acceptTM_run, Turing.FinTM, Turing.FinTM.Computes, Turing.FinTM.ComputesInTime, Turing.FinTM.ComputesInTime.output_unique, Turing.FinTM.computesInTime_iff, Turing.MultiTapeTM.initCfg, Turing.MultiTapeTM.runFrom}
\difficulty{5}
Let $M$ be a finite-state binary-alphabet Turing machine and
$p : \{0,1\}^{*} \to \{0,1\}$ a predicate such that on every input $x$ the
machine $M$ eventually halts with the single bit $p(x)$ as its output. Let
$M^{\mathrm{acc}}$ be the halting recognizer of $M$, which simulates $M$ with
output suppressed while recording the last bit $M$ emits (initially $0$), halts
when $M$ halts with recorded bit $1$, and otherwise enters a loop state in which
it does nothing forever. Then for every input $x$, $M^{\mathrm{acc}}$ halts on
$x$ (for some output $w$ and some $t$, it halts within $t$ steps with output $w$)
if and only if $p(x) = 1$.
\end{lemma}

\begin{definition}[Machine prepending a fixed word]
\lean{Complexity.prefixTM}
\label{Complexity.prefixTM}
\leanok
\uses{Turing.FinTM}
For a binary string $w = w_0 w_1 \cdots w_{m-1}$, the prefixing machine $P_w$ is
the binary-alphabet Turing machine with no work tapes, states
$\{0, 1, \ldots, m\}$ and start state $0$. In state $i < m$ it emits $w_i$ and
moves to state $i+1$ without moving its input head; in state $m$, if the input
head scans a bit $b$ it emits $b$, moves the input head right and stays in state
$m$, and if it scans a blank it halts without emitting anything.
\end{definition}

\begin{definition}[Configurations of the prefixing machine]
\lean{Complexity.prefixCfg}
\label{Complexity.prefixCfg}
\leanok
\uses{Turing.Cfg}
For binary strings $w$ and $x$, an optional state
$q \in \{0, 1, \ldots, |w|\} \cup \{\bot\}$ ($\bot$ meaning halted), an
input-head position $p \in \{0, 1, \ldots, |x|+1\}$ (positions $0$ and $|x|+1$
being the blank cells flanking $x$) and an output string $u$, this is the
configuration of a machine with no work tapes and state set
$\{0, 1, \ldots, |w|\}$ on input $x$ that is in state $q$, has its input head at
position $p$, and has $u$ on its output tape.
\end{definition}

\begin{lemma}[Emission phase of the prefixing machine]
\lean{Complexity.prefixTM_emit}
\label{Complexity.prefixTM_emit}
\leanok
\uses{Complexity.prefixCfg, Complexity.prefixTM, Turing.Action.apply, Turing.Cfg.ext_zero_tapes, Turing.MultiTapeTM.initCfg, Turing.MultiTapeTM.runFrom, Turing.MultiTapeTM.runFrom_succ_eq_step', Turing.MultiTapeTM.step}
\difficulty{4}
Let $w = w_0 \cdots w_{m-1}$ and $x$ be binary strings, and let $P_w$ be the
machine with no work tapes and states $\{0, \ldots, m\}$ which, starting in
state $0$, emits $w_i$ and moves to state $i+1$ in each state $i < m$, and in
state $m$ copies the input bit by bit, halting at the first blank. For every
$i \le m$, after $i$ steps on input $x$ the machine $P_w$ is in state $i$, its
input head is still on the first input cell, and its output tape holds the first
$i$ bits $w_0 \cdots w_{i-1}$ of $w$.
\end{lemma}

\begin{lemma}[Copy phase of the prefixing machine]
\lean{Complexity.prefixTM_copy}
\label{Complexity.prefixTM_copy}
\leanok
\uses{Complexity.prefixCfg, Complexity.prefixTM, Turing.Action.apply, Turing.Cfg.ext_zero_tapes, Turing.Cfg.inputSymbol, Turing.MultiTapeTM.runFrom, Turing.MultiTapeTM.runFrom_succ_eq_step', Turing.inputSymbolInner, Turing.moveInputPos, Turing.moveInputPos_pos_of_ne_right}
\difficulty{4}
Let $w = w_0 \cdots w_{m-1}$ and $x = x_1 \cdots x_n$ be binary strings, and let
$P_w$ be the machine with no work tapes and states $\{0, \ldots, m\}$ which emits
$w_i$ and advances to state $i+1$ in each state $i < m$, and in state $m$ emits
the scanned input bit and moves its input head right, halting at a blank. For
every $i \le n$, running $P_w$ on input $x$ for $i$ steps from the configuration
in state $m$ with the input head on the first input cell and output $w$ yields
the configuration in state $m$ with the input head on cell $i+1$ and output
$w\,x_1 \cdots x_i$.
\end{lemma}

\begin{lemma}[The prefixing machine computes $x \mapsto wx$ in linear time]
\lean{Complexity.prefixTM_computes}
\label{Complexity.prefixTM_computes}
\leanok
\uses{Complexity.prefixCfg, Complexity.prefixTM, Complexity.prefixTM_copy, Complexity.prefixTM_emit, Turing.Action.apply, Turing.Cfg.inputSymbol, Turing.FinTM.ComputesFunInTime, Turing.FinTM.computesInTime_iff, Turing.MultiTapeTM.runFrom_add, Turing.MultiTapeTM.runFrom_succ_eq_step', Turing.MultiTapeTM.step}
\difficulty{4}
Let $w = w_0 \cdots w_{m-1}$ be a binary string and $P_w$ the machine with no
work tapes and states $\{0, \ldots, m\}$ which first emits $w_0, \ldots, w_{m-1}$
one per step without moving its input head, then copies the input bit by bit
while moving right, halting at the first blank. Then for every binary input $x$,
$P_w$ halts within $|w| + |x| + 1$ steps with output the concatenation $wx$.
\end{lemma}

\begin{lemma}[Running time of fixed-code pairing]
\lean{Complexity.fixedPair_computes}
\label{Complexity.fixedPair_computes}
\leanok
\uses{Complexity.prefixTM, Complexity.prefixTM_computes, Turing.FinTM.ComputesFunInTime, Turing.pairEncode}
\difficulty{4}
Let $\alpha = \alpha_1 \cdots \alpha_m$ be a binary string, and for a binary
string $x$ let $\langle \alpha, x \rangle = \alpha_1\alpha_1 \cdots
\alpha_m\alpha_m\,0\,1\,x$ be the pairing encoding ($\alpha$ with every bit
doubled, then the separator $01$, then $x$). The machine with no work tapes that
first emits the fixed word $\alpha_1\alpha_1 \cdots \alpha_m\alpha_m\,01$ one bit
per step and then copies its input bit by bit, halting at the first blank,
halts on every input $x$ within $2|\alpha| + |x| + 3$ steps with output
$\langle \alpha, x \rangle$.
\end{lemma}

\begin{lemma}[Fixed-code pairing is polynomial-time computable]
\lean{Complexity.fixedPair_polyTime}
\label{Complexity.fixedPair_polyTime}
\leanok
\uses{Complexity.PolyTimeComputable, Complexity.fixedPair_computes, Complexity.prefixTM, Turing.FinTM.ComputesInTime.mono, Turing.pairEncode}
\difficulty{3}
Let $\alpha$ be a fixed binary string, and for binary strings $x$ let
$\langle \alpha, x \rangle$ be the pairing encoding: $\alpha$ with every bit
doubled, followed by $01$, followed by $x$. Then the map
$x \mapsto \langle \alpha, x \rangle$ is polynomial-time computable: some
finite-state binary-alphabet Turing machine computes it within
$C \cdot (n+1)^c$ steps on inputs of length $n$, for constants
$C, c \in \mathbb{N}$.
\end{lemma}

\begin{theorem}[The halting language is $\mathsf{NP}$-hard]
\lean{Complexity.HALT_NPHard}
\label{Complexity.HALT_NPHard}
\leanok
\uses{Complexity.HALT, Complexity.HALT_pairEncode_eq_true_iff, Complexity.NPHard, Complexity.NP_subset_EXP, Complexity.acceptTM, Complexity.acceptTM_halts_iff, Complexity.fixedPair_polyTime, Turing.FinTM.ComputesFunInTime, Turing.FinTM.ComputesFunInTime.computes, Turing.FinTM.one_work_tape_binary, Turing.MachineCode, Turing.MachineCode.decode_encode, Turing.MultiTapeTM.indicator, Turing.exists_codeTM, Turing.pairEncode}
\difficulty{5}
Let $c$ be any representation scheme for Turing machines: a total decoding of
bit strings into coded one-work-tape machines together with an encoding, such
that a machine's code followed by any number of $1$s decodes back to that
machine. Let $\mathrm{HALT}_c$ be the set of bit strings of the form
$\langle \alpha, x \rangle$ ($\alpha$ with every bit doubled, then $01$, then
$x$) such that the machine decoded from $\alpha$ halts on input $x$. Then
$\mathrm{HALT}_c$ is $\mathsf{NP}$-hard: every language in $\mathsf{NP}$ is
polynomial-time Karp reducible to $\mathrm{HALT}_c$ [AB09, Exercise~2.8].
\end{theorem}

\begin{theorem}[The halting language is not in $\mathsf{NP}$]
\lean{Complexity.HALT_not_mem_NP}
\label{Complexity.HALT_not_mem_NP}
\leanok
\uses{Complexity.HALT, Complexity.HALT_not_computable, Complexity.NP, Complexity.NP_subset_EXP, Turing.EffectiveMachineCode, Turing.FinTM.ComputesFunInTime, Turing.FinTM.ComputesFunInTime.computes, Turing.FinTM.DecidesInTime, Turing.MultiTapeTM.indicator}
\difficulty{4}
Let $c$ be an effective representation scheme for Turing machines: a total
decoding of bit strings into coded machines with an encoding that decoding
inverts (also after padding the code with $1$s), equipped with a finite-state
binary-alphabet machine that, within some time bound, computes from every bit
string the canonical serialization of the machine it decodes to. Let
$\mathrm{HALT}_c$ be the set of bit strings of the form $\langle \alpha, x \rangle$
($\alpha$ with every bit doubled, then $01$, then $x$) such that the machine
decoded from $\alpha$ halts on input $x$. Then $\mathrm{HALT}_c \notin
\mathsf{NP}$ [AB09, Exercise~2.8].
\end{theorem}
